\scanpage{138}
% lines 1–2
Soit $f:X\unsure{-|\to} S$ un morphisme se factorisant en $X\overset{\varphi}{-}\Sigma\overset{\psi}{-}S$, avec $\psi$ fini, et $\varphi$ tel que $\underline O_\Sigma\overset{\sim}{\longrightarrow}\varphi_{*}(\underline O_X)$.

% lines 3–8
1) Alors \unsure{on a}
\[
\unsure{\cdot|0}\to\mathrm{Pic}(\Sigma)-\mathrm{Pic}(X)-\mathrm{Pic}'(X/S)
\]
[suite exacte provenant de la suite spectrale de Leray].

\deleted{\unsure{et}} \uline{Si} $X/\Sigma$ admet une section, le dernier hom est un isom, i.e. dans ce cas
\[
\mathrm{Pic}'(X/S)\simeq\mathrm{Pic}(\deleted{\Sigma}X)/\mathrm{Pic}(\deleted{\unsure{X|\Sigma}}\Sigma)
\]
% lines 9–16
Ce dernier fait se voit ainsi: \deleted{Pour} Soit $g:\Sigma\longrightarrow X$ une section de $X/\Sigma$ ; alors \inserted{ii) \unsure{puis|pour}} \deleted{\unsure{tout}} \unsure{1|Les} Module inversible $g$-rigidifié sur $X$ n'a pas d'aut. non trivial, \inserted{iii)} tout Module inversible $\underline L$ sur $X$ est \deleted{\unsure{de isom|de façon}} $\mathrm{Pic}'(X/S)$ \uline{équivalent} à \deleted{l'image} un Module inversible $g$-rigidifié \inserted{[savoir $\underline L\otimes\varphi^{*}(g^{*}(\underline L))$]} [déterminé à isom unique près].

[On utilise le fait que $\mathrm{Pic}'(\Sigma/S)=0$].

% last line
On suppose pour la suite \deleted{\unsure{f}} que, quel \unsure{de|est} que soit
% archive sheet number: 137

